% Options for packages loaded elsewhere
\PassOptionsToPackage{unicode}{hyperref}
\PassOptionsToPackage{hyphens}{url}
\PassOptionsToPackage{dvipsnames,svgnames,x11names}{xcolor}
%
\documentclass[
]{article}
\usepackage{amsmath,amssymb}
\usepackage{iftex}
\ifPDFTeX
  \usepackage[T1]{fontenc}
  \usepackage[utf8]{inputenc}
  \usepackage{textcomp} % provide euro and other symbols
\else % if luatex or xetex
  \usepackage{unicode-math} % this also loads fontspec
  \defaultfontfeatures{Scale=MatchLowercase}
  \defaultfontfeatures[\rmfamily]{Ligatures=TeX,Scale=1}
\fi
\usepackage{lmodern}
\ifPDFTeX\else
  % xetex/luatex font selection
  \setmainfont[]{DejaVu Serif}
  \setmonofont[]{DejaVu Sans Mono}
\fi
% Use upquote if available, for straight quotes in verbatim environments
\IfFileExists{upquote.sty}{\usepackage{upquote}}{}
\IfFileExists{microtype.sty}{% use microtype if available
  \usepackage[]{microtype}
  \UseMicrotypeSet[protrusion]{basicmath} % disable protrusion for tt fonts
}{}
\makeatletter
\@ifundefined{KOMAClassName}{% if non-KOMA class
  \IfFileExists{parskip.sty}{%
    \usepackage{parskip}
  }{% else
    \setlength{\parindent}{0pt}
    \setlength{\parskip}{6pt plus 2pt minus 1pt}}
}{% if KOMA class
  \KOMAoptions{parskip=half}}
\makeatother
\usepackage{xcolor}
\usepackage[margin=1in]{geometry}
\usepackage{longtable,booktabs,array}
\usepackage{calc} % for calculating minipage widths
% Correct order of tables after \paragraph or \subparagraph
\usepackage{etoolbox}
\makeatletter
\patchcmd\longtable{\par}{\if@noskipsec\mbox{}\fi\par}{}{}
\makeatother
% Allow footnotes in longtable head/foot
\IfFileExists{footnotehyper.sty}{\usepackage{footnotehyper}}{\usepackage{footnote}}
\makesavenoteenv{longtable}
\setlength{\emergencystretch}{3em} % prevent overfull lines
\providecommand{\tightlist}{%
  \setlength{\itemsep}{0pt}\setlength{\parskip}{0pt}}
\setcounter{secnumdepth}{-\maxdimen} % remove section numbering
\usepackage{newunicodechar}
\newunicodechar{★}{\ensuremath{\star}}
\newunicodechar{✓}{\ensuremath{\checkmark}}
\newunicodechar{√}{\ensuremath{\surd}}
\newunicodechar{≈}{\ensuremath{\approx}}
\ifLuaTeX
  \usepackage{selnolig}  % disable illegal ligatures
\fi
\IfFileExists{bookmark.sty}{\usepackage{bookmark}}{\usepackage{hyperref}}
\IfFileExists{xurl.sty}{\usepackage{xurl}}{} % add URL line breaks if available
\urlstyle{same}
\hypersetup{
  colorlinks=true,
  linkcolor={Maroon},
  filecolor={Maroon},
  citecolor={Blue},
  urlcolor={Blue},
  pdfcreator={LaTeX via pandoc}}

\author{}
\date{}

\begin{document}

\hypertarget{chapter-2-softmax-and-temperature}{%
\section{Chapter 2 --- Softmax and
temperature}\label{chapter-2-softmax-and-temperature}}

\begin{verbatim}
Type: theory + notebook
Ordering: theory-first
Depends on: 1 (scaled dot-product attention)
\end{verbatim}

\hypertarget{motivation}{%
\subsection{2.1 Motivation}\label{motivation}}

Attention scores \(s_{ij}\) are arbitrary reals. To turn them into
usable attention weights we need a map \(\mathbb{R}^n \to \Delta^{n-1}\)
(the probability simplex) that is non-negative, sum-to-one,
differentiable, monotone in each coordinate, and --- ideally --- the
\emph{exponential-family} canonical link. Softmax satisfies all of these
and is the unique such map up to scaling. Its shape parameter,
\emph{temperature}, controls how selectively attention concentrates.
Understanding softmax deeply is understanding attention's
information-retrieval dynamics.

\hypertarget{storyline}{%
\subsection{2.2 Storyline}\label{storyline}}

Softmax does one job with four useful consequences.

\textbf{The job.} Given logits \(s \in \mathbb{R}^n\), softmax emits the
Gibbs distribution \(p_i = e^{s_i/T} / Z\) on \(\{1, \dots, n\}\), with
temperature \(T > 0\) controlling concentration.

\textbf{Useful consequence 1: shift invariance.} Adding a constant to
every logit leaves the distribution unchanged. This is why the standard
numerical implementation subtracts \(\max_i s_i\) before exponentiating
--- it prevents overflow without affecting the answer.

\textbf{Useful consequence 2: smoothness.} Softmax is \(C^\infty\), with
a clean Jacobian
\(\partial p_i / \partial s_j = p_i(\delta_{ij} - p_j) / T\). This
enables gradient flow to the upstream \(Q, K\) projections. Its failure
at saturation (Lemma 3.3 preview) is the central reason scaling matters.

\textbf{Useful consequence 3: temperature controls entropy.} As
\(T \to 0^+\), softmax approaches a one-hot at \(\arg\max s\); as
\(T \to \infty\), it approaches uniform. The entropy of \(p\) is a
smooth monotone function of \(T\) for fixed \(s\). This is the
mechanism's information-control knob.

\textbf{Useful consequence 4: log-sum-exp as a soft max.}
\(\text{LSE}(s) = \log \sum_i e^{s_i}\) is a smooth upper bound on
\(\max_i s_i\), and \(\nabla \text{LSE} = \text{softmax}\). This is why
``softmax'' is called that --- it is the gradient of a smooth-max
function.

\textbf{Temperature in attention.} In scaled dot-product attention the
\(1/\sqrt{d_k}\) factor is exactly an effective temperature:
\(\text{softmax}(QK^\top / \sqrt{d_k}) = \text{softmax}_{\sqrt{d_k}}(QK^\top)\).
Chapter 3 derives why this particular temperature is the right one;
Chapter 2 is the tool-building that makes Chapter 3 possible.

\textbf{Regime notes.} Softmax itself is stateless --- no train / val /
inference asymmetry in the operator. Temperature, however, can differ
across regimes: - Training uses whatever temperature is built into the
architecture (typically \(\sqrt{d_k}\) for attention, \(1\) for the
output layer). - \textbf{Inference-time sampling} from a language model
often uses a \emph{different} temperature (sampling temperature) on the
output softmax --- concentrating or broadening the next-token
distribution independently of how the model was trained. This is a
\emph{post-training} knob. - Masking (causal or attention masks) is
applied additively to the \emph{pre-softmax} scores, usually by adding
\(-\infty\) to disallowed positions, producing zero weight there after
softmax.

\hypertarget{rigorous-core}{%
\subsection{2.3 Rigorous core}\label{rigorous-core}}

\textbf{Definition 2.1 (Softmax with temperature).} For
\(s \in \mathbb{R}^n\) and \(T > 0\),
\[\text{softmax}_T(s)_i = \frac{e^{s_i / T}}{\sum_{j=1}^n e^{s_j / T}}.\]
We write \(\text{softmax}(s) := \text{softmax}_1(s)\) for the
unit-temperature default.

\textbf{Proposition 2.2 (Shift invariance).} For any
\(c \in \mathbb{R}\),
\(\text{softmax}_T(s + c \mathbf{1}) = \text{softmax}_T(s)\). (The
constant \(e^{c/T}\) factors out of numerator and denominator.)

\textbf{Proposition 2.3 (Jacobian).} Let \(p = \text{softmax}(s)\). Then
\[\frac{\partial p_i}{\partial s_j} = p_i(\delta_{ij} - p_j).\] For
\(\text{softmax}_T\), the Jacobian is
\(\tfrac{1}{T} p_i(\delta_{ij} - p_j)\). \emph{Quotient rule, case-split
on \(i = j\).}

\textbf{Corollary 2.3.1 (Mass conservation).} The rows and columns of
the Jacobian each sum to zero:
\(\sum_i \partial p_i / \partial s_j = 0\). \emph{Why:} softmax maps to
the simplex, so \(\sum_i p_i = 1\) is constant; any perturbation in
\(s_j\) preserves this.

\textbf{Proposition 2.4 (Entropy bounds).} For
\(p = \text{softmax}(s)\),
\(H(p) = -\sum_i p_i \log p_i \in [0, \log n]\). The upper bound is
attained iff all \(s_i\) are equal; the lower bound is approached as the
pre-softmax margin grows without bound. \emph{Upper: Jensen / Lagrange.
Lower: entropy non-negative on any distribution.}

\textbf{Proposition 2.5 (Saturation margin).} For
\(s = (0, \Delta, 0, \dots, 0) \in \mathbb{R}^n\) (winner at index 2
with margin \(\Delta\)), the winning probability is
\[p_2 = \frac{e^\Delta}{e^\Delta + (n-1)} = \frac{1}{1 + (n-1) e^{-\Delta}}.\]
To achieve \(p_2 > 1 - \epsilon\) requires
\(\Delta > \log((n-1)/\epsilon)\) to leading order.

\textbf{Corollary 2.5.1.} Margin requirements are \emph{logarithmic} in
\(n\) --- doubling the number of keys costs one additional nat of
margin. This is why attention over long contexts is not hopeless despite
the \(n\)-way competition; the cost to distinguish a clear winner from
\(n\) distractors grows only as \(\log n\).

\textbf{Proposition 2.6 (Log-sum-exp).} Let
\(\text{LSE}(s) = \log \sum_i e^{s_i}\). Then:

\begin{enumerate}
\def\labelenumi{(\alph{enumi})}
\tightlist
\item
  \(\nabla \text{LSE}(s) = \text{softmax}(s)\);
\item
  \(\max_i s_i \le \text{LSE}(s) \le \max_i s_i + \log n\);
\item
  \(\text{LSE}\) is convex.
\end{enumerate}

\emph{Sketches: (a) direct \(\partial_{s_i} \text{LSE} = p_i\). (b)
sandwich \(e^m \le \sum_i e^{s_i} \le n e^m\). (c) Hessian is the
softmax Jacobian, which is PSD as the covariance of the categorical
\(p\) --- Exercise 2.1.}

\textbf{Proposition 2.7 (Temperature as inverse logit scaling).}
\(\text{softmax}_T(s) = \text{softmax}(s/T)\). Therefore the
\(1/\sqrt{d_k}\) factor in attention corresponds to effective
temperature \(T_{\text{eff}} = \sqrt{d_k}\).

\hypertarget{numerical-example-shift-invariance-and-saturation-by-hand}{%
\subsection{2.4 Numerical example --- shift invariance and saturation by
hand}\label{numerical-example-shift-invariance-and-saturation-by-hand}}

\textbf{Shift invariance.} \(\text{softmax}((1, 2, 3))\)
vs.~\(\text{softmax}((101, 102, 103))\): identical. The second form
overflows a naive implementation; subtracting \(103\) first recovers the
first form. Hand-computed:
\[\text{softmax}((1, 2, 3)) = \frac{1}{e + e^2 + e^3} (e, e^2, e^3) \approx (0.090, 0.245, 0.665).\]

\textbf{Saturation by margin.} For \(n = 10\) and
\(\Delta \in \{1, 3, 5, 10\}\) the winning probability
\(p_{\text{win}} = e^\Delta / (e^\Delta + 9)\) is:

\begin{longtable}[]{@{}ll@{}}
\toprule\noalign{}
\(\Delta\) & \(p_{\text{win}}\) \\
\midrule\noalign{}
\endhead
\bottomrule\noalign{}
\endlastfoot
1 & \(\approx 0.252\) \\
3 & \(\approx 0.691\) \\
5 & \(\approx 0.943\) \\
10 & \(\approx 0.9996\) \\
\end{longtable}

Margin \(\Delta \approx \log(n / \epsilon)\) is where \(p_{\text{win}}\)
crosses \(1 - \epsilon\): for \(n = 10\), \(\epsilon = 0.01\), that's
\(\Delta \approx 6.9\) --- the table above brackets it.

\textbf{LSE vs.~max.} For the same vector \((0, 5, 0, 0, \dots, 0)\)
with \(n = 10\), \(\max = 5\),
\(\text{LSE} = 5 + \log(1 + 9e^{-5}) \approx 5.061\). Tight at large
margin, off by up to \(\log n \approx 2.3\) at zero margin.

\hypertarget{mechanics-check}{%
\subsection{2.5 Mechanics check}\label{mechanics-check}}

\begin{itemize}
\tightlist
\item
  \textbf{Exponential} --- converts any real vector to positive entries,
  enforcing non-negativity.
\item
  \textbf{Normalizer \(Z = \sum_j e^{s_j}\)} --- enforces sum-to-one,
  turning a positive-valued map into a probability distribution.
\item
  \textbf{Shift-invariance} --- makes the implementation numerically
  stable (subtract the max), and means softmax only cares about score
  \emph{differences}, not absolute scales.
\item
  \textbf{Temperature \(T\) (or equivalently, pre-softmax scaling)} ---
  the \emph{only} knob controlling concentration. Every selectivity
  argument in attention is ultimately a temperature argument.
\item
  \textbf{Not hard-max} --- softmax's smoothness is what makes attention
  trainable (Prop 1.4 + Lemma 3.3). Argmax would give identical
  forward-pass semantics in the \(T \to 0\) limit but zero gradient.
\item
  \textbf{Asymptotics at \(T \to 0\) and \(T \to \infty\)} --- one-hot
  and uniform respectively. The entire useful range of attention lives
  between these two extremes.
\end{itemize}

\hypertarget{exercises}{%
\subsection{2.6 Exercises}\label{exercises}}

\textbf{2.1 ★★} \emph{{[}Mechanism dissection --- variance of the
categorical.{]}} Let \(p = \text{softmax}(s)\) and let \(X\) be the
categorical random variable with \(P(X = i) = p_i\). For any function
\(f: \{1, \dots, n\} \to \mathbb{R}\), compute \(\text{Var}_p(f(X))\) in
terms of \(p\) and \(f\). Show that the softmax Jacobian (Prop 2.3) is
precisely the covariance matrix of the indicator vector \(\mathbf{1}_X\)
of this categorical. What does this say about the Jacobian's
eigenvalues?

\textbf{2.2 ★★} \emph{{[}Failure-mode construction --- temperature
pathologies.{]}} Construct pre-softmax inputs
\(s^{(1)}, s^{(2)} \in \mathbb{R}^n\) such that:

\begin{enumerate}
\def\labelenumi{(\alph{enumi})}
\item
  at \(T = 1\), \(\text{softmax}(s^{(1)})\) and
  \(\text{softmax}(s^{(2)})\) have the same top-1 probability, but at
  \(T = 10\) they differ in top-1 by more than \(0.1\);
\item
  the entropy of \(\text{softmax}_T(s)\) is non-monotone in \(T\) over
  some range.
\end{enumerate}

For (b), state whether non-monotonicity is even possible. Prove or
disprove.

\textbf{2.3 ★★★} \emph{{[}Theorem about the optimization --- LSE
duality.{]}} Consider the variational characterization
\[\text{LSE}(s) = \max_{p \in \Delta^{n-1}} \left\{ \langle p, s \rangle + H(p) \right\},\]
where \(\Delta^{n-1}\) is the probability simplex and
\(H(p) = -\sum_i p_i \log p_i\).

\begin{enumerate}
\def\labelenumi{(\alph{enumi})}
\item
  Prove the identity: show the maximizer is \(p = \text{softmax}(s)\)
  and the maximum value is \(\text{LSE}(s)\).
\item
  Generalize to temperature: express \(T \cdot \text{LSE}(s/T)\) as a
  similar variational problem with the entropy term scaled
  appropriately. What does temperature correspond to in this dual
  picture?
\item
  Interpret the duality in attention: what is the
  ``entropy-regularized'' optimization that softmax attention is
  implicitly solving at each query position?
\end{enumerate}

\textbf{2.4 ★★★} \emph{{[}Failure mode --- hard attention and gradient
estimators.{]}} A hard-attention layer selects one key per query via
argmax instead of softmax. Gradients vanish almost everywhere (Prop 1.4
failure).

\begin{enumerate}
\def\labelenumi{(\alph{enumi})}
\item
  Sketch the Gumbel-softmax estimator: add Gumbel noise to the scores,
  softmax at temperature \(T\), anneal \(T\) toward zero during
  training. Why does this preserve gradients while approaching hard
  selection?
\item
  Sketch the straight-through estimator: forward pass is argmax;
  backward pass pretends the argmax was the identity map (or, more
  commonly, pretends it was softmax). Why is this \emph{biased} in the
  gradient sense, and what is the bias? Name at least one failure mode
  in terms of training dynamics.
\item
  Compare: under what circumstances would you choose Gumbel-softmax over
  straight-through, and vice versa? What role does the scale of \(n\)
  (number of keys) play?
\end{enumerate}

\textbf{2.5 ★★★} \emph{{[}Cross-chapter combination ---
predict-then-verify.{]}} The notebook will compute softmax entropy of
\(\text{softmax}(\mathbf{s})\) for
\(\mathbf{s} \sim \mathcal{N}(0, \sigma^2 I_n)\) as a function of
\((\sigma, n)\). Before running, predict:

\begin{enumerate}
\def\labelenumi{(\alph{enumi})}
\item
  for fixed \(n\), entropy as a function of \(\sigma\) --- direction,
  asymptotes at \(\sigma \to 0\) and \(\sigma \to \infty\), inflection
  behavior;
\item
  for fixed \(\sigma\), entropy as a function of \(n\) --- direction,
  scaling with \(n\);
\item
  is there a combined \((\sigma, n)\)-quantity that controls entropy?
  Conjecture one, defend it, and state what you expect to see in a
  heatmap of entropy over \((\sigma, n)\). \emph{Hint: the score
  variance of a single element is \(\sigma^2\); relate this to Prop
  2.5.}
\end{enumerate}

State confidence per part with a one-line reason.

\hypertarget{before-the-notebook-02_softmax_temperature.ipynb}{%
\subsection{\texorpdfstring{2.7 Before the notebook
(\texttt{02\_softmax\_temperature.ipynb})}{2.7 Before the notebook (02\_softmax\_temperature.ipynb)}}\label{before-the-notebook-02_softmax_temperature.ipynb}}

Write down the following predictions \emph{before} opening the notebook.

\begin{enumerate}
\def\labelenumi{\arabic{enumi}.}
\tightlist
\item
  \textbf{Shift invariance check.} Direction and expected numerical
  error of the difference
  \(\text{softmax}((1, 2, 3)) - \text{softmax}((101, 102, 103))\) under
  the safe (max-subtracting) and naive (no subtraction) implementations.
\item
  \textbf{Entropy vs.~temperature.} Fix
  \(s \sim \mathcal{N}(0, I_{50})\). Plot of \(H(\text{softmax}_T(s))\)
  as \(T\) ranges over \([0.01, 100]\) on a log-x axis. State the two
  asymptotes and the direction of the curve between them.
\item
  \textbf{Saturation margin.} For \(n = 10\) and target
  \(p_{\text{win}} = 0.99\), predict \(\Delta^\star\) and the
  closed-form dependence on \(n\) (from Corollary 2.5.1).
\item
  \textbf{Entropy heatmap over \((\sigma, n)\).} From Exercise 2.5(c):
  what combined quantity do you conjecture controls entropy? Predict the
  level-set shape of the entropy heatmap over
  \(\log \sigma \times \log n\) --- are level sets vertical, horizontal,
  or diagonal lines?
\item
  \textbf{Jacobian PSD check.} Softmax Jacobian is a covariance matrix
  (Exercise 2.1). Predict: is it PSD or strictly PD? What is its rank?
\item
  \textbf{LSE sandwich.} For \(s \sim \mathcal{N}(0, I_n)\), predict the
  gap \(\text{LSE}(s) - \max_i s_i\) as a function of \(n\). Does it
  grow like \(\log n\) (the worst-case bound) or something smaller?
\end{enumerate}

For each: direction, magnitude, confidence. The notebook will confirm or
update each prediction.

\end{document}
